%
%	Einfuehrung
%

\pagebreak
\section{Introduction}

\onehalfspacing

\subsection{NIS2 and CIS Benchmarks}

NIS2 entered into force on 16 January 2023, and the Member States now have 21 months, until 17 October 2024, to transpose its measures into national law.\footnote{See \textit{Negreiro-Achiaga, M. (2023)}: The NIS2 Directive. \cite{nisBrief}}

\subsection{Kubernetes}

Kubernetes 1.30

\begin{figure}[H]
\centering
\caption {Kubernetes 1.30 Release Logo}
\includegraphics[width=\linewidth]{images/k8s-1.30.png}
\label{fig:uwubernetes}
\end{figure}

\subsection{Gender-neutral Pronouns}

Our society is becoming more open, inclusive, and gender-fluid, and now I think it's time to think about using gender-neutral pronouns in scientific texts, too. Two well-known researchers, Abigail C. Saguy and Juliet A. Williams, both from UCLA, propose to use singular they/them instead: "The universal singular they is inclusive of people who identify as male, female or nonbinary."\footnote{\textit{Saguy, A. (2020)}: Why We Should All Use They/Them Pronouns. \cite{pronouns}} The aim is to support an inclusive approach in science through gender-neutral language. 

In this paper, I'll attempt to follow this suggestion and invite all my readers to do the same for future articles. Thank you!

If you're not sure about the definitions of gender and sex and how to use them, have a look at the definitions\footnote{See \textit{APA (2021)}: Definitions Related to Sexual Orientation. \cite{apaDefinitions}} by the American Psychological Association.

\subsection{Online Learning and Loneliness}

According to a recent study conducted by the American Psychological Association, social distancing over an extended period can increase loneliness and significantly affect people's health.\footnote{See \textit{Luchetti, M. (2020)}: The trajectory of loneliness in response to COVID-19. \cite{apaLoneliness}}

\subsection{Methodology \& Research Question}

Peerforming an Experiment.\footnote{See \textit{Genau, L. (2020)}: Ein Experiment in deiner Abschlussarbeit durchführen. \cite{expScribbr}}
